
%
% This document has designed to be read through as a PDF first.
% Click on the ``Typeset'' button in the toolbar to generate the PDF
% if it is not already visible.
%

%%% PREAMBLE %%%
% You probably want to skip to \begin{document} if this is your first time.

\documentclass[article,oneside]{memoir}
% This command goes at the beginning of every document
% [oneside,article] are two of many options that can be chosen
%   oneside makes each page have the same layout, for printing on only one side 
%     of the paper (change it to twoside to see the difference)
%   article means we're writing a short document only and won't be using special 
%     chapter headings
%   a4paper changes the page dimensions for A4 sized paper 

%%% PACKAGES %%%

\usepackage{graphicx} % Add graphics capabilities
\usepackage{float}
\usepackage{booktabs} % ``Proper'' table layout
\usepackage{amsmath}  % Better maths support
\usepackage[colorlinks=true,linkcolor=red]{hyperref} % Hyperlink capabilities

\usepackage{memhfixc} % This package is required to resolve incompatibilities
                      % with the memoir class & the hyperref package
                      
 \usepackage[T1]{fontenc}
 
                      
                      
%%\usepackage{pdfsync}
%%\usepackage{/Library/texmf/tex/latex/pdfsync}
% This package is used to tell TeXShop where things are in the PDF file.
% Command-click at any spot in the PDF and it will jump to the corresponding
% location in the source file.

%%% COMMANDS %%%

% These are two examples of how you can define your own commands.
% These two are used to save on typing.
\newcommand{\pkg}[1]{\textsf{#1}}
\newcommand{\lshort}{\href{http://www.ctan.org/tex-archive/info/lshort/english/lshort.pdf}{The Not So Short Introduction to \LaTeX{}}}

\usepackage{titlesec}
    \titleformat{\chapter}[display]
        {\normalfont\Large\bfseries}    % Make title \large and bold
        { Chapter~\thechapter}     % Write Chapter X in \huge font
        {20pt}                          % 20pt distance between Chapter X and title
        {}
        [\vspace{1ex}\titlerule]        % Add a line below chapter title
   % {\huge Chapter~\thechapter}     % Write Chapter X in \huge font     
        
% Define the title, author and date of the document.
% If the date is undefined, the current date is substituted.
\title {Synthetic Aperture Radar Point Target Response}
\author{Sasan Ardalan Extra Class Amateur Radio, Ph.D.}
\date{}



%%% BODY OF THE DOCUMENT: %%%
\begin{document}


\thispagestyle{empty}
\maketitle
\thispagestyle{empty} % Removes the page number from the title page

% Generates a title based on the \title, \author, and \date commands in the preamble.
\begin{abstract}
 
The algorithms for computing the  point target response in Synthetic Aperture Radar (SAR) will be presented. The target modeling and simulations will be performed following the procedure developed by McDonough, et. al. (1985) [1] for SEASAT. Simulation results will be provided through the block diagram modeling   of the SAR system with Capsim\textsuperscript{\textregistered}[2]. The SAR project has been a part of the Capsim\textsuperscript{\textregistered} distribution since 1990. The research on SAR was conducted  by the author while a Professor at NC State University in 1987. The GitHub repository was created in November 2025.

\end{abstract}
\paragraph{\copyright{Copyright  1987-2006 Sasan Ardalan}}



Permission is granted to copy, distribute and/or modify this
document under the terms of the GNU Free Documentation License,
Version 1.2 or any later version published by the Free Software
Foundation; with no Invariant Sections, no Front-Cover Texts, and
no Back-Cover Texts. A copy of the license is included in the
section entitled "GNU Free Documentation License".


\pagebreak

 
 \setcounter{tocdepth}{2}
 \tableofcontents
 \pagebreak
 
\listoffigures


\pagebreak

% Please add the following required packages to your document preamble:
% \usepackage{booktabs}






\chapter{Introduction}

The algorithms for computing the  point target response in Synthetic Aperture Radar (SAR) will be presented. The target modeling and simulations will be performed following the procedure developed by McDonough, et. al. (1985) [1] for SEASAT. The algorithms and simulation results for a point target in Synthetic Aperture Radar (SAR) will be presented. The  typical SEASAT-SAR geometry is shown  in Figure \ref{fig:geometry} and with parameters given in Table 1. A technique by Fitch [3]  will also be used to avoid a Hilbert transform. 




 \begin{figure}[H]
  \centering 
    \includegraphics[width=4in]{./SAR_Geometry_2.pdf}
  \caption{SAR Geometry}
\label{fig:geometry}
\end{figure}



\chapter{Point Target Response}

\begin{table}[]
\centering
\caption{SAR Parameters}
\label{my-label}
\begin{tabular}{@{}lllll@{}}
\toprule
Item & Parameter         & Value  & Units  &  \\ \midrule
1    & Carrier Frequency & 1275   & MHz    &  \\
2    & Pulse Chirp Rate  & 0.5621 & MHz/uS &  \\
3    & Pulse Duration    & 33.8   & uS     &  \\
4	&Pulse Bandwidth 	&19.0	&MHz 	& \\
5	&Center Frequency (IF) 	&11.38	&MHz & \\
6   &Pulse Repetition Rate(PRF)   &1645.0	&Hz   & \\
7   &Sampling Rate &45.03	&MHz	& \\
8   &Doppler Frequency  &1150.0	   &Hz &   \\
9   &Doppler Rate of Change   &501.27    &Hz  & \\
10  &Platform Velocity   &7.0	&Km/s & \\
11  &Integration Time &2.0   &s   & \\
12  &Azimuth Sample Time &0.0	&s & \\
13 &tc	& 1.0	&s & \\
14 &Point Range, rp	&840.0  &kM &\\
15  &Point Azimuth tp in index units &1000 &index &\\


\end{tabular}
\end{table}
Consider a point target arbitrarily placed at position  ($r_{p},\tau_{p}$). The radar will fly at some height $h$, and velocity $v$. The slant range to the target changes as a function of viewing position as shown by
\begin{equation}
\label{eq:slantRange}
r(t,\tau) = r_p+\frac{v^2}{2r_p}(\tau_c+\tau_p-\tau)^2
\end{equation}
where $\tau$ is the time variable in the azimuth direction, $\tau_c$ coincides with the mid point of the integration time, and $\tau_p$ is the azimuthal time location of the point target relative to $\tau_c$. 
To derive \ref{eq:slantRange} , we note that 

\begin{equation}
\label{eq:slantRange2}
r^2(t,\tau) = r_p^2+\left (v\times(\tau_c+\tau_p-\tau)\right )^2
\end{equation}

Then using a Taylor series expansion,

\begin{equation}
\label{eq:slantRange3}
r(t,\tau) = r_p +\frac{v^2}{2r_p}(\tau_c+\tau_p-\tau)^2  - \frac{v^4}{2r_p^3}\left (\tau_c+\tau_p-\tau\right )^4 + ...
\end{equation}


Note that  $\frac{v\times(\tau_c+\tau_p-\tau)}{r_P} \ll 1 $ so we get Equation  \ref{eq:slantRange} .

Equation \ref{eq:slantRange} shows that the range history of the target during the integration time for forming an aperture is described by a parabola. For example, if we consider $\tau_p$ to coincide with $\tau_c$, then the slant range will be a maximum when the target first enters the beam at $\tau=0$. It will then reach a minimum when $\tau = \tau_c$, and will again increase to some maximum as the target exits the beam. Now, if the $\tau_p$ is not identically located at $\tau_c$, then the beginning and ending slant ranges will not be equal. The difference in the two range positions is known as range walk, and together with the parabolic curvature introduced into the range, the problem is known as range migration.




To obtain high resolution in the range direction, the radar transmits an FM chirped pulse of the form,

\begin{equation}
\label{ }
s(t)=\cos 2 \pi (f_0 t +\pi K_r t^2) ~~~ |t| \leq  \Delta t/2
\end{equation}

where $s(t)$ is the transmitted signal, $f_0$ is the carrier frequency of the radar, $K_r$ is the chirp rate and $\Delta t$ is the pulse width.




Each pulse has a duration of $33.8~\mu  $sec and a chirp rate of 0.5621~MHz$/\mu $sec or a pulse bandwidth $B_r$ of 19 MHz which results in a slant range resolution $\delta r= \frac{c}{2}B_r \sin\theta $ of 25 meters. 
\par

\par

Figure \ref{fig:chirp} shows plots of the time domain and frequency domain signal for the radar FM chirped pulse for range.
\par

The received pulse $g(t)$ will be delayed by the round trip time of $ 2 \frac{r(\tau )}{c}$, or,

\begin{equation}
\label{ }
g(t)=\cos \left [ \left (2 \pi f_0 (t-2r(\tau)/c \right ) + \pi K_r \left ( t -2 r(\tau)/c \right )^2 \right  ]~~| t-2r(\tau)/c | \leq \Delta t/2
\end{equation}

where $r(\tau)$ is the slant range from the radar at time $\tau$, to the target at position $(r_p,\tau_p)$, and $c$ is the speed of light. To minimize the required sampling rate, the signal is then mixed with a local oscillator at the carrier frequency 1275 MHz and down-converted to some intermediate frequency of $f_1 $= 11.38 MHz resulting in,

\begin{equation}
\label{eq:matchedFilter}
g(t)=\cos \left [ \left (2 \pi f_1 t-2r(\tau)/\lambda  \right )+ \pi K_r \left ( t -2 r(\tau)/c \right )^2 \right ] 
\end{equation}

\section{Range Compression}

To get any useful information regarding the position of the target, the down-converted received signal can be correlated with a reference function. Thus, when the two signals are coincident at some  time and have the same shape the output of the correlation is at a maximum. The correlation is performed as the convolution of $g'(t)$ and the replica $s'(-t)$ as [1],


\begin{equation}
\label{eq:convolution}
s'(-t)*g'(t) = \Delta t/2 Re \left [ \sin \left (\pi B_r \left ( t-2r(\tau)/c \right ) \right ) /(\pi B_r \left ( t-2r(\tau)/c \right ) 
\times e^{-j4 \pi r(\tau)/\lambda } e^{j2 \pi f_1 t } \right ]  
\end{equation}


where
\begin{equation}
\label{eq:convolution-condition}
| t-\frac{2r(\tau)}{c} | \leq \Delta t
\end{equation}



and  $B_r$ is the bandwidth.



Hence, the time domain convolution in the above expression results in the desired range compressed signal. Now, the reference function for the pulse compression has $U=\tau f_s$ or 33.8 $\mu sec \times 45.03 ~MHz =1530$ points. Thus the time domain correlation will require on the order of $U\times U $ or $2\times 10^6$ multiplications. The number of operations, however, can be significantly reduced to $2U\log U $ or $9.7\times 10^3$ if the time domain convolution is performed as multiplication in the frequency domain. The range compression in the frequency domain is obtained from:





\begin{enumerate}
\item Take the DFFT ($g'(t)) $ to get $G'(f)$ of length $N$.
\item Take the DFFT ($s'(t)$) to get $S'(f)$ of length $N$.
\item Truncate the spectrums to half of the original lengths.
\item Range compressed = $FFT^{-1} \left (  G'(f) S'(f) \right )$ of length $\frac{N}{2}$.
\item $(r_0)$ position of point target $= \max ( | FFT^{-1} \left ( G'(f) S'(f) \right ) | ) $.
\end{enumerate}

At the end of operation 4 we are left with the right hand side of  Equation \ref{eq:convolution}. Step number 3 is necessary since the azimuth compression requires a complex sequence for correlation. This important approach for real to complex conversion using FFTs was described by Fitch in 1988 [3] and it bypasses the need for a Hilbert transform. Basically, the complex sequence has a real part equal to samples of the real signal taken at half the original rate and a spectrum equal to the first half of the original real sequence spectrum. The real to complex conversion is accomplished by using the first half of the spectral coefficients.

\par

The result of the matched filter for range compression is illustrated in the plot in Figure \ref{fig:range_compression}.

\par 


The range compression operation begins by fixing an arbitrary reference point within the matrix with a known $r(t_{ref},\tau_{ref} )$ and computing its impulse response $g(t)$ using Equation \ref{eq:matchedFilter}. Then, the response $g(t)$ to the fixed value of $\tau$, which is the azimuthal position of $r(t,\tau)$ of Equation \ref{eq:slantRange}. The time extents of both $s(t)$ and $g(t)$ are $33.8~\mu \sec $ and the signals are sampled at $f_s$ or 45.03 MHz. Each pulse is then compressed using a 2048 point FFT according to the procedure described above.


\par

The result of compressing a single pulse is shown in Figure \ref{fig:range_compression}. It basically demonstrates the result of the correlation output from Equation  \ref{eq:convolution} showing an envelope of $\frac{\sin x}{x}$ which is maximized when $t=2r_P/c$, where $r_p$ is the range to the point target. Hence, the range location, $r$, of the point target has been recovered. 
Figure \ref{fig:range_compression_sketch} illustrates the FM Chirped Pulse and the Compression and Resolution Related to Bandwidth. To show  that the range compression is indeed a $\frac{\sin(x)}{x}$ function, Figure \ref{fig:range_compression_sinx_over_x}  shows the range compression of Figure \ref{fig:range_compression} zoomed in.

\par


The range compression process is repeated by moving in the azimuthal direction, fixing the slant range to the point target $r(\tau)$ and transmitting pulses at the PRF rate.  For this particular simulation, using  an integration time of 2 seconds and a PRF of 1645 Hz, a total of 3290 pulses were transmitted.
\par




If we consider a number of successive echoes, due to the large change in the range from the satellite to the target during integration time $T_i$,  the range compressed point response will be a parabolic curve as shown by Equation \ref{eq:slantRange} and :

\begin{equation}
\label{ }
\frac{Z^2}{r_0 c} - \frac{\tau}{2} < Y < \frac{Z^2}{r_0 c} +  \frac{\tau}{2} 
\end{equation}

and

\begin{equation}
\label{ }
- \frac{vT_i}{2} < Z < \frac{vT_i}{2}
\end{equation}

where $Y$ is the start of each pulse, $Z=v\tau$, and $\tau'$ is the pulse duration.



\begin{figure}[H]
  \centering 
    \includegraphics[height=8cm]{SAR_Image_Formation/FFT_FM_Chirp_SAR.png}
  \caption{SAR Range FM Chirp Radar Signal}
\label{fig:chirp}
\end{figure}

\begin{figure}[H]
  \centering 
    \includegraphics[height=8cm]{SAR_Image_Formation/range_compression_plot.png}
  \caption{SAR Image Range Compression}
\label{fig:range_compression}
\end{figure}


\begin{figure}[H]
  \centering 
    \includegraphics[height=8cm]{SAR_Image_Formation/chirp_pulse_compression.pdf}
  \caption{SAR Image Range FM Chirped Radar and Compression }
\label{fig:range_compression_sketch}
\end{figure}

\begin{figure}[H]
  \centering 
    \includegraphics[height=8cm]{SAR_Image_Formation/range_compression_plot_zoomed_bar_chart.png}
  \caption{SAR Image Range Compression Zoomed Discrete Showing Sin(x)/x Response}
\label{fig:range_compression_sinx_over_x}
\end{figure}

As the range returns are received a matrix is formed as shown in Figure \ref{fig:sar_image_record}. 

 \begin{figure}[H]
  \centering 
    \includegraphics[width=4in]{./SAR_Record_Radar_Returns.png}
  \caption{SAR Record Radar Returns}
\label{fig:sar_image_record}
\end{figure}




\section{Azimuth Compression}

Now, the complex exponential term $e^{-j4\pi r(\tau)\lambda}$ appearing in Equation \ref{eq:convolution}, represents the azimuth phase history which is encoded in the range compressed signal. If $r(\tau)$ is expanded about $\tau_c$,
\begin{equation}
\label{ }
r(\tau) =r(\tau_c) + \dot{r}(\tau_c)(\tau -\tau_c)+\frac{1}{2} \ddot{r} (\tau_c)(\tau -\tau_c)^2+...
\end{equation}
where
\begin{equation}
\label{ }
f_{DC} = -\frac{2\dot{r}}{\lambda}
\end{equation}
is the Doppler frequency and, 

\begin{equation}
\label{ }
K_{az} = -\frac{2\ddot{r}}{\lambda}
\end{equation}
is the rate of change of Doppler frequency.

$\dot{r}$ and $\ddot{r}$ are the velocity and acceleration vectors respectively.
\par

Re-writting $r(\tau) $ in terms of $f_{DC}$ and $K_{az}$ we get [1]:

\begin{equation}
\label{ }
\begin{split}
g(\tau) = e^{-j4\pi r(\tau)/\lambda }= e^{-j4\pi r_c/\lambda } e^{j2\pi f_{DC}(\tau -\tau_c) + j\pi K_{az} (\tau -\tau_c)^2} \\
~~|\tau -\tau_c| \leq \Delta \tau/2 ~~ and  ~~ 0 ~ elsewhere
\end{split}
\end{equation}

The matched filter for the azimuth compression will have an impulse response of the form:

\begin{equation}
\label{eq:azimuthMatchedFilter}
h(\tau)= e^{j2\pi (f_{DC} \tau -0.5 K_{az} \tau^2) } ~~|\tau| \leq \Delta \tau/2
\end{equation}

\par

Figure \ref{fig:sar_azimuth_ref} shows a plot of the Azimuth complex matched filter $h(\tau)$.

\par

For the azimuth compression an operation similar to the range compression is required. The range compressed lines are read back into memory from external disk storage, and the data are corner turned such that consecutive azimuth cells are aligned for a given range line. In actuality, at this stage a range walk and curvature correction is required since the range data are collected along a parabolic curve in the azimuth direction. Each range compressed line will then have a $g(t)$ which will then be convolved with $h(\tau)$ to get[1]:


\begin{equation}
\label{ }
\begin{split}
b(\tau)=g(\tau)* h(\tau) = e^{j2\pi f_{DC}\tau)} 
\times   e^{-j2\pi (2r_c/\lambda + f_{DC} \tau_c)}     \times \\
\Delta  \tau \left ( \pi B_D (\tau -\tau_c) \right ) /(\pi B_D (\tau -\tau_c)) ~~ | \tau -\tau_c| \leq \Delta \tau /2
\end{split}
\end{equation}



where $B_D= 2 f_{DC} $. The matched filter response $h(\tau)$, is generated with a $\tau$ spacing of $\frac{1}{PRF}$ and $f_{DC} $= 1150 Hz and $ K_{az}$= 501.3  Hz. Once again because of computational efficiency the above convolution is performed in the frequency domain:

\begin{enumerate}
\item  Take the $DFFT(g(\tau)) $ to get $G(f)$.
\item  Take the $DFFT(h(\tau)) $ to get $H(f)$.
\item  Azimuth compressed $b(\tau) = FFT^{-1} \left ( G(f)H(f) \right ) $
\item   $\tau $ position of point target$= max \left  ( | b(\tau ) | \right )$
\end{enumerate}



We note that since both $g(\tau)$ and $h(\tau)$ are complex, the original length of the sequence for the FFT's were used. The above process is repeated until each of the range compressed lines are compressed in the azimuth. The output of the correlator $\frac{\sin \pi B_D(\tau -\tau_c)}{\pi B_D(\tau -\tau_c}$ is a maximum when $\tau =\tau_c$, or when the azimuth position of the point target is aligned with the mid point of the beam. 



\begin{figure}[H]
  \centering 
    \includegraphics[height=8cm]{SAR_Image_Formation/SAR_Azimuth_Reference.png}
  \caption{SAR Azimuth Range Reference (Complex)}
\label{fig:sar_azimuth_ref}
\end{figure}



\chapter{SAR Image for Point Targets and Block Diagram Modeling and Simulation}

The Capsim\textsuperscript{\textregistered}  Block Diagram of the SAR Image Formation is shown in Figure \ref{fig:sar_block_diagram}. The complete C code is provided in the GitHub  Repository [2].
\par
An illustration of the SAR image prior to range  compression and azimuth compression for a point target created by the block diagram modeling and simulation  tool is shown in Figure \ref{fig:create_image}.

\par



\begin{figure}[H]
  \centering 
    \includegraphics[width=4in]{./SAR_Block_Diagram_Capsim.png}
  \caption{Capsim\textsuperscript{\textregistered} SAR Block Diagram}
\label{fig:sar_block_diagram}
\end{figure}




\begin{figure}[H]
  \centering 
    \includegraphics[height=10cm]{SAR_Image_Formation/SAR_Image_Point_Target_grs.jpg}
  \caption{SAR Image of Point Target Range 851 km Azimuth Index 350 Range Time 67.6 $\mu$s}
\label{fig:create_image}
\end{figure}


The SAR image for  point targets are shown in Figures \ref{fig:sar_image_formation}, and \ref{fig:sar_image_formation_2}.




\begin{figure}[H]
  \centering 
    \includegraphics[height=10cm]{SAR_Image_Formation/sar_image_formation_840p8_400.jpg}
  \caption{SAR Image Formation Point Target Range 840.8 km  Azimuth Index 400 Range Time 33.8 $\mu$s}
\label{fig:sar_image_formation}
\end{figure}




\begin{figure}[H]
  \centering 
    \includegraphics[height=10cm]{SAR_Image_Formation/sar_image_formation_842p5_700.jpg}
  \caption{SAR Image Formation Point Target  Range 842.5 km Azimuth Index 700 Range Time 33.8 $\mu$s}
\label{fig:sar_image_formation_2}
\end{figure}






\chapter{Auto-Focusing}

The fundamental parameters in SAR target response are the Doppler frequency and the rate of change of the frequency. These parameters are solely dependent on the rate of change of the relative motion characteristics between the SAR sensor and the position of the target. Therefore, a precise knowledge of the spacecraft state vector is necessary in order to determine these parameters. In situations where a precise satellite tracking and ephemeris are not given and the Doppler spectrum and frequency sweep rate are not known apriori, these parameters must be derived from the radar data. Estimating the Doppler parameters is the topic of autofocus and clutter-locking. This adds to the system complexity, and a controller must be designed to not only generate the reference functions for the azimuth, but also to  provide a feedback loop to refine and readjust the Doppler parameters.


\chapter{SAR Block Diagram Simulation Repository on GitHub}
Complete C code for  SAR signal processing is available for the Capsim\textsuperscript{\textregistered} Block Diagram Simulation and Modeling Environment. Visit the GitHub repository at:
\par
https://github.com/silicondsp/synthetic-aperture-radar
\par 
\par 
The repository includes all the C code blocks  and documentation.
 Also visit
\par
https://www.ccdsp.org
\par

\pagebreak
\chapter{References}

1- Robert N. McDonough, Barry E. Raff and Joyce L. Kerr, Image formation from spaceborne synthetic aperture radar signals, 1985 APL Technical Digest  Vol. 6, No. 4 pp 300-312
\subsection{  }
2- Capsim\textsuperscript{\textregistered} Block Diagram Modeling and Simulation Tool, Silicon DSP Corporation, https://www.silicondsp.com, 1989-2025. Also see the repository on GitHub: https://github.com/silicondsp/capsim-tmk
\subsection{  }
3- Fitch, J.P. , Synthetic Aperture Radar. 1988 Springer-Verlag New York Berlin Heidelbergh.





\pagebreak
\chapter{GNU Free Documentation License}
GNU Free Documentation License
Version 1.3, 3 November 2008
\par
\noindent
Copyright \copyright  2000, 2001, 2002, 2007, 2008 Free Software Foundation, Inc.  <http://fsf.org/>
Everyone is permitted to copy and distribute verbatim copies of this license document, but changing it is not allowed.

\subsection{0. PREAMBLE}


The purpose of this License is to make a manual, textbook, or other functional and useful document "free" in the sense of freedom: to assure everyone the effective freedom to copy and redistribute it, with or without modifying it, either commercially or noncommercially. Secondarily, this License preserves for the author and publisher a way to get credit for their work, while not being considered responsible for modifications made by others.\\
This License is a kind of "copyleft", which means that derivative works of the document must themselves be free in the same sense. It complements the GNU General Public License, which is a copyleft license designed for free software.\\
We have designed this License in order to use it for manuals for free software, because free software needs free documentation: a free program should come with manuals providing the same freedoms that the software does. But this License is not limited to software manuals; it can be used for any textual work, regardless of subject matter or whether it is published as a printed book. We recommend this License principally for works whose purpose is instruction or reference.\\
\subsection{1. APPLICABILITY AND DEFINITIONS}


This License applies to any manual or other work, in any medium, that contains a notice placed by the copyright holder saying it can be distributed under the terms of this License. Such a notice grants a world-wide, royalty-free license, unlimited in duration, to use that work under the conditions stated herein. The "Document", below, refers to any such manual or work. Any member of the public is a licensee, and is addressed as "you". You accept the license if you copy, modify or distribute the work in a way requiring permission under copyright law.
A "Modified Version" of the Document means any work containing the Document or a portion of it, either copied verbatim, or with modifications and/or translated into another language.\\
A "Secondary Section" is a named appendix or a front-matter section of the Document that deals exclusively with the relationship of the publishers or authors of the Document to the Document's overall subject (or to related matters) and contains nothing that could fall directly within that overall subject. (Thus, if the Document is in part a textbook of mathematics, a Secondary Section may not explain any mathematics.) The relationship could be a matter of historical connection with the subject or with related matters, or of legal, commercial, philosophical, ethical or political position regarding them.\\
The "Invariant Sections" are certain Secondary Sections whose titles are designated, as being those of Invariant Sections, in the notice that says that the Document is released under this License. If a section does not fit the above definition of Secondary then it is not allowed to be designated as Invariant. The Document may contain zero Invariant Sections. If the Document does not identify any Invariant Sections then there are none.
The "Cover Texts" are certain short passages of text that are listed, as Front-Cover Texts or Back-Cover Texts, in the notice that says that the Document is released under this License. A Front-Cover Text may be at most 5 words, and a Back-Cover Text may be at most 25 words.\\
A "Transparent" copy of the Document means a machine-readable copy, represented in a format whose specification is available to the general public, that is suitable for revising the document straightforwardly with generic text editors or (for images composed of pixels) generic paint programs or (for drawings) some widely available drawing editor, and that is suitable for input to text formatters or for automatic translation to a variety of formats suitable for input to text formatters. A copy made in an otherwise Transparent file format whose markup, or absence of markup, has been arranged to thwart or discourage subsequent modification by readers is not Transparent. An image format is not Transparent if used for any substantial amount of text. A copy that is not "Transparent" is called "Opaque".\\
Examples of suitable formats for Transparent copies include plain ASCII without markup, Texinfo input format, LaTeX input format, SGML or XML using a publicly available DTD, and standard-conforming simple HTML, PostScript or PDF designed for human modification. Examples of transparent image formats include PNG, XCF and JPG. Opaque formats include proprietary formats that can be read and edited only by proprietary word processors, SGML or XML for which the DTD and/or processing tools are not generally available, and the machine-generated HTML, PostScript or PDF produced by some word processors for output purposes only.\\
The "Title Page" means, for a printed book, the title page itself, plus such following pages as are needed to hold, legibly, the material this License requires to appear in the title page. For works in formats which do not have any title page as such, "Title Page" means the text near the most prominent appearance of the work's title, preceding the beginning of the body of the text.\\
The "publisher" means any person or entity that distributes copies of the Document to the public.\\
A section "Entitled XYZ" means a named subunit of the Document whose title either is precisely XYZ or contains XYZ in parentheses following text that translates XYZ in another language. (Here XYZ stands for a specific section name mentioned below, such as "Acknowledgements", "Dedications", "Endorsements", or "History".) To "Preserve the Title" of such a section when you modify the Document means that it remains a section "Entitled XYZ" according to this definition.\\
The Document may include Warranty Disclaimers next to the notice which states that this License applies to the Document. These Warranty Disclaimers are considered to be included by reference in this License, but only as regards disclaiming warranties: any other implication that these Warranty Disclaimers may have is void and has no effect on the meaning of this License.

\subsection{2. VERBATIM COPYING}



You may copy and distribute the Document in any medium, either commercially or noncommercially, provided that this License, the copyright notices, and the license notice saying this License applies to the Document are reproduced in all copies, and that you add no other conditions whatsoever to those of this License. You may not use technical measures to obstruct or control the reading or further copying of the copies you make or distribute. However, you may accept compensation in exchange for copies. If you distribute a large enough number of copies you must also follow the conditions in section 3.\\
You may also lend copies, under the same conditions stated above, and you may publicly display copies.

\subsection{3. COPYING IN QUANTITY}



If you publish printed copies (or copies in media that commonly have printed covers) of the Document, numbering more than 100, and the Document's license notice requires Cover Texts, you must enclose the copies in covers that carry, clearly and legibly, all these Cover Texts: Front-Cover Texts on the front cover, and Back-Cover Texts on the back cover. Both covers must also clearly and legibly identify you as the publisher of these copies. The front cover must present the full title with all words of the title equally prominent and visible. You may add other material on the covers in addition. Copying with changes limited to the covers, as long as they preserve the title of the Document and satisfy these conditions, can be treated as verbatim copying in other respects.\\
If the required texts for either cover are too voluminous to fit legibly, you should put the first ones listed (as many as fit reasonably) on the actual cover, and continue the rest onto adjacent pages.\\
If you publish or distribute Opaque copies of the Document numbering more than 100, you must either include a machine-readable Transparent copy along with each Opaque copy, or state in or with each Opaque copy a computer-network location from which the general network-using public has access to download using public-standard network protocols a complete Transparent copy of the Document, free of added material. If you use the latter option, you must take reasonably prudent steps, when you begin distribution of Opaque copies in quantity, to ensure that this Transparent copy will remain thus accessible at the stated location until at least one year after the last time you distribute an Opaque copy (directly or through your agents or retailers) of that edition to the public.\\
It is requested, but not required, that you contact the authors of the Document well before redistributing any large number of copies, to give them a chance to provide you with an updated version of the Document.

\subsection{4. MODIFICATIONS}




You may copy and distribute a Modified Version of the Document under the conditions of sections 2 and 3 above, provided that you release the Modified Version under precisely this License, with the Modified Version filling the role of the Document, thus licensing distribution and modification of the Modified Version to whoever possesses a copy of it. In addition, you must do these things in the Modified Version:\\
$\bullet$	A. Use in the Title Page (and on the covers, if any) a title distinct from that of the Document, and from those of previous versions (which should, if there were any, be listed in the History section of the Document). You may use the same title as a previous version if the original publisher of that version gives permission. \\
$\bullet$	B. List on the Title Page, as authors, one or more persons or entities responsible for authorship of the modifications in the Modified Version, together with at least five of the principal authors of the Document (all of its principal authors, if it has fewer than five), unless they release you from this requirement. \\
$\bullet$	C. State on the Title page the name of the publisher of the Modified Version, as the publisher. \\
$\bullet$	D. Preserve all the copyright notices of the Document. \\
$\bullet$	E. Add an appropriate copyright notice for your modifications adjacent to the other copyright notices. \\
$\bullet$	F. Include, immediately after the copyright notices, a license notice giving the public permission to use the Modified Version under the terms of this License, in the form shown in the Addendum below. \\
$\bullet$	G. Preserve in that license notice the full lists of Invariant Sections and required Cover Texts given in the Document's license notice. \\
$\bullet$	H. Include an unaltered copy of this License. \\
$\bullet$	I. Preserve the section Entitled "History", Preserve its Title, and add to it an item stating at least the title, year, new authors, and publisher of the Modified Version as given on the Title Page. If there is no section Entitled "History" in the Document, create one stating the title, year, authors, and publisher of the Document as given on its Title Page, then add an item describing the Modified Version as stated in the previous sentence. \\
$\bullet$	J. Preserve the network location, if any, given in the Document for public access to a Transparent copy of the Document, and likewise the network locations given in the Document for previous versions it was based on. These may be placed in the "History" section. You may omit a network location for a work that was published at least four years before the Document itself, or if the original publisher of the version it refers to gives permission. \\
$\bullet$	K. For any section Entitled "Acknowledgements" or "Dedications", Preserve the Title of the section, and preserve in the section all the substance and tone of each of the contributor acknowledgements and/or dedications given therein. \\
$\bullet$	L. Preserve all the Invariant Sections of the Document, unaltered in their text and in their titles. Section numbers or the equivalent are not considered part of the section titles. \\
$\bullet$	M. Delete any section Entitled "Endorsements". Such a section may not be included in the Modified Version. \\
$\bullet$	N. Do not retitle any existing section to be Entitled "Endorsements" or to conflict in title with any Invariant Section. \\
$\bullet$	O. Preserve any Warranty Disclaimers. \\
If the Modified Version includes new front-matter sections or appendices that qualify as Secondary Sections and contain no material copied from the Document, you may at your option designate some or all of these sections as invariant. To do this, add their titles to the list of Invariant Sections in the Modified Version's license notice. These titles must be distinct from any other section titles.\\
You may add a section Entitled "Endorsements", provided it contains nothing but endorsements of your Modified Version by various parties?for example, statements of peer review or that the text has been approved by an organization as the authoritative definition of a standard.
You may add a passage of up to five words as a Front-Cover Text, and a passage of up to 25 words as a Back-Cover Text, to the end of the list of Cover Texts in the Modified Version. Only one passage of Front-Cover Text and one of Back-Cover Text may be added by (or through arrangements made by) any one entity. If the Document already includes a cover text for the same cover, previously added by you or by arrangement made by the same entity you are acting on behalf of, you may not add another; but you may replace the old one, on explicit permission from the previous publisher that added the old one.
The author(s) and publisher(s) of the Document do not by this License give permission to use their names for publicity for or to assert or imply endorsement of any Modified Version.

\subsection{5. COMBINING DOCUMENTS}



You may combine the Document with other documents released under this License, under the terms defined in section 4 above for modified versions, provided that you include in the combination all of the Invariant Sections of all of the original documents, unmodified, and list them all as Invariant Sections of your combined work in its license notice, and that you preserve all their Warranty Disclaimers.\\
The combined work need only contain one copy of this License, and multiple identical Invariant Sections may be replaced with a single copy. If there are multiple Invariant Sections with the same name but different contents, make the title of each such section unique by adding at the end of it, in parentheses, the name of the original author or publisher of that section if known, or else a unique number. Make the same adjustment to the section titles in the list of Invariant Sections in the license notice of the combined work.
In the combination, you must combine any sections Entitled "History" in the various original documents, forming one section Entitled "History"; likewise combine any sections Entitled "Acknowledgements", and any sections Entitled "Dedications". You must delete all sections Entitled "Endorsements".

\subsection{6. COLLECTIONS OF DOCUMENTS}


You may make a collection consisting of the Document and other documents released under this License, and replace the individual copies of this License in the various documents with a single copy that is included in the collection, provided that you follow the rules of this License for verbatim copying of each of the documents in all other respects.
You may extract a single document from such a collection, and distribute it individually under this License, provided you insert a copy of this License into the extracted document, and follow this License in all other respects regarding verbatim copying of that document.

\subsection{7. AGGREGATION WITH INDEPENDENT WORKS}


A compilation of the Document or its derivatives with other separate and independent documents or works, in or on a volume of a storage or distribution medium, is called an "aggregate" if the copyright resulting from the compilation is not used to limit the legal rights of the compilation's users beyond what the individual works permit. When the Document is included in an aggregate, this License does not apply to the other works in the aggregate which are not themselves derivative works of the Document.
If the Cover Text requirement of section 3 is applicable to these copies of the Document, then if the Document is less than one half of the entire aggregate, the Document's Cover Texts may be placed on covers that bracket the Document within the aggregate, or the electronic equivalent of covers if the Document is in electronic form. Otherwise they must appear on printed covers that bracket the whole aggregate.

\subsection{8. TRANSLATION}


Translation is considered a kind of modification, so you may distribute translations of the Document under the terms of section 4. Replacing Invariant Sections with translations requires special permission from their copyright holders, but you may include translations of some or all Invariant Sections in addition to the original versions of these Invariant Sections. You may include a translation of this License, and all the license notices in the Document, and any Warranty Disclaimers, provided that you also include the original English version of this License and the original versions of those notices and disclaimers. In case of a disagreement between the translation and the original version of this License or a notice or disclaimer, the original version will prevail.
If a section in the Document is Entitled "Acknowledgements", "Dedications", or "History", the requirement (section 4) to Preserve its Title (section 1) will typically require changing the actual title.

\subsection{9. TERMINATION}


You may not copy, modify, sublicense, or distribute the Document except as expressly provided under this License. Any attempt otherwise to copy, modify, sublicense, or distribute it is void, and will automatically terminate your rights under this License.
However, if you cease all violation of this License, then your license from a particular copyright holder is reinstated (a) provisionally, unless and until the copyright holder explicitly and finally terminates your license, and (b) permanently, if the copyright holder fails to notify you of the violation by some reasonable means prior to 60 days after the cessation.
Moreover, your license from a particular copyright holder is reinstated permanently if the copyright holder notifies you of the violation by some reasonable means, this is the first time you have received notice of violation of this License (for any work) from that copyright holder, and you cure the violation prior to 30 days after your receipt of the notice.
Termination of your rights under this section does not terminate the licenses of parties who have received copies or rights from you under this License. If your rights have been terminated and not permanently reinstated, receipt of a copy of some or all of the same material does not give you any rights to use it.

\subsection{10. FUTURE REVISIONS OF THIS LICENSE}


The Free Software Foundation may publish new, revised versions of the GNU Free Documentation License from time to time. Such new versions will be similar in spirit to the present version, but may differ in detail to address new problems or concerns. See http://www.gnu.org/copyleft/.
Each version of the License is given a distinguishing version number. If the Document specifies that a particular numbered version of this License "or any later version" applies to it, you have the option of following the terms and conditions either of that specified version or of any later version that has been published (not as a draft) by the Free Software Foundation. If the Document does not specify a version number of this License, you may choose any version ever published (not as a draft) by the Free Software Foundation. If the Document specifies that a proxy can decide which future versions of this License can be used, that proxy's public statement of acceptance of a version permanently authorizes you to choose that version for the Document.

\subsection{11. RELICENSING}


?Massive Multiauthor Collaboration Site" (or "MMC Site") means any World Wide Web server that publishes copyrightable works and also provides prominent facilities for anybody to edit those works. A public wiki that anybody can edit is an example of such a server. A "Massive Multiauthor Collaboration" (or "MMC") contained in the site means any set of copyrightable works thus published on the MMC site.
"CC-BY-SA" means the Creative Commons Attribution-Share Alike 3.0 license published by Creative Commons Corporation, a not-for-profit corporation with a principal place of business in San Francisco, California, as well as future copyleft versions of that license published by that same organization.
"Incorporate" means to publish or republish a Document, in whole or in part, as part of another Document.
An MMC is "eligible for relicensing" if it is licensed under this License, and if all works that were first published under this License somewhere other than this MMC, and subsequently incorporated in whole or in part into the MMC, (1) had no cover texts or invariant sections, and (2) were thus incorporated prior to November 1, 2008.
The operator of an MMC Site may republish an MMC contained in the site under CC-BY-SA on the same site at any time before August 1, 2009, provided the MMC is eligible for relicensing.
\end{document}

